% !TeX root = ../ba-lentz.tex

\documentclass[ba-lentz.tex]{subfile}

\begin{document}

{\let\clearpage\relax \chapter{Introduction: Hybrid approaches to seasonal drought prediction}}\label{introduction}

\noindent Extreme droughts have a great impact on agricultural, social, and economic conditions.
They occur in almost all climate zones and vary widely in duration, intensity, and regional extent.
A distinction can be made between meteorological, agricultural, and socioeconomic droughts.
While metereological droughts are defined as the unusual and extreme absence of sufficient precipitation, agricultural and socioeconomic definitions focus on the societal impacts of metereological droughts.
However, the prediction of meteorological drought events even with state-of-the-art global climate models (GCMs) often lacks significant prediction skill.
On seasonal time scales, GCMs predict precipitation and temperatures, which are consequently used to calculate drought indices, such as the Standardised Precipitation Evaporation Index (SPEI).
Despite huge advancements in seasonal climate predictions in the last decades, climate models still possess significant biases and uncertainties when predicting drought events \parencite{wu_dynamic-lstm_2022}.
Especially after the loss of initialization effects, GCMs still show significant errors in their dynamical representation of metereological droughts \parencite{wu_dynamic-lstm_2022}.
In alternative approaches, statistical methods can be used to directly predict the occurence and intensity of drought events from historical data.
Statistical methods have different limitations.
For example, basic statistical methods are often limited in their capability to represent nonlinear dynamics reliably \parencite{hao_seasonal_2018}.
More advanced statistical methods like machine learning (ML), while being able to learn and reproduce non-linear dynamics from their training data, lack interpretability and a good representation of overall distributions. \\
In order to combine the strengths of both methods, hybrid approaches based on a mix of GCM prediction and ML postprocessing have recently gotten a lot of attention \parencite{li_post-processing_2021,wu_dynamic-lstm_2022}.
Due to the sparse amount of training data and a focus on reliability and explainability, these hybrid approaches often use relatively simple ML models.
In comparison to other fields in metereological or climate research, modern ML techniques, like deep neural networks or transformer models, seem to be less suited than reliable and cost-efficient solutions like random forest models \parencite{thomasson_improving_2023,feng_bias_2024,karsisto_post-processing_2025}.
Therefore, this work will focus on the use of decision tree based ML methods for the postprocessing of drought predictions.
Especially extreme gradient boosting (XGBoost) has shown to be a reliable, efficient and flexible option for forecast postprocessing \parencite{feng_bias_2024}.
However, XGBoost also exibhits struggles with extreme outliers and to correct for systematic biases in the variance and distribution of its input data \parencite{feng_bias_2024}.
Thus, an intermediate statistical step will be performed.
In this step, quantile mapping (QM) matches the original hindcasts distribution to that of the reference data.
Since QM cannot by design correct the trend or time evolution of its input data, the combination of QM and XGBoost should be able to account for each other's weaknesses.
Thus, I hope thath this combined framework significantly improves the original hindcast's prediction skill.\\
Despite the success of ML postprocessing in some fields of climate science, they have not yet been successfully applied to seasonal predictions.
Especially the prediction of extreme events, such as droughts, has only recently gotten more attention in seasonal prediction systems.
Nonetheless, there is still significant room for improvement, especially in the postprocessing of drought predictions.
The DWD for example does not possess an operational seasonal forecast postprocessing including ML techniques to this day.
This work tries to close this gap.
By combining dynamical predictions of climatic variables on seasonal timescales with a statistical postprocessing step, I will analyse if ML models can be used in order to improve seasonal drought prediction. \\
There are two aspects of drought prediction that can be improved.
Either one can try to improve the mean prediction of a chosen drought indicator.
Alternatively, the prediction of extreme drought events can be improved.
While this work will focus on the deterministic mean improvement, it will also evaluate, if the ML postprocessing has an impact on the extreme event prediction.
Often the SPEI is used to measure the severity of droughts.
However, the standardization process as well as its requirement of full values instead of anomalies make it a more complicated postprocessing target.
Therefore, the climatic water balance (CWB) will be chosen as a drought indicator.
The CWB includes both precipitation and potential evapotranspiration (PET) in its calculation and the basis from which the SPEI is calculated.
Without the standardization and the possibility to calculate monthly anomalies, the postprocessing of the CWB with ML methods can additionally hold a much larger potential for generalization.
Consequently, the method developed in this work also holds importance for other potential use cases, such as the postprocessing of other climate variables or on different timescales. \\
In order outline the experimental setup, first an overview of the utilized methods and data will be given in chapter \ref{data_methods}.
Both the original dynamical hindcasts and the observational data used for the ML training will be descibed in detail in section \ref{data}.
The used ML framework as well as additional post- and preprocessing steps will be documented in \ref{methods}.
Afterwards, the results of the subsequent analysis will be shown in chapter \ref{results}.
At last, the results will be discussed and a conclusion will be drawn in section \ref{conclusion} in regard to this hybrid approach's ability to significantly improve the seasonal prediction of drought events.



\end{document}